\documentclass[10pt, twocolumn, landscape, a4paper]{article}
% CEPRU | 3er Examen - Ciclo Ordinario 2017-I / Tema P | Grupo A..D
\usepackage[margin=1.5cm]{geometry}
\usepackage[utf8]{inputenc}
\usepackage[spanish,es-noshorthands]{babel}
\usepackage{amssymb, amsmath, amsbsy}
\usepackage{tasks}
\usepackage{enumerate}
\usepackage{enumitem}
\usepackage{graphicx} % \scalebox{1.4}{$$}
\usepackage[pdftex]{hyperref}
\hypersetup{pdftitle={},pdfauthor={Luis Carrillo Gutiérrez}}

\fontfamily{cmss}\selectfont
\renewcommand{\familydefault}{\sfdefault}
\pagestyle{empty} 
\setlength{\columnsep}{1.2in}

\begin{document}
\subsubsection*{Álgebra}
\begin{itemize}
\item{El dominio de la relación de variable real : \\ $R = \big\{ (x,y) \in \mathbb{R}^2\;/\;xy^2 - x^2y^2 - x = 2 \big\}$ es : 
\begin{tasks}(2)
\task{$\big< -\!\infty,\;- 2\big]  \cup \big<0, 1\big>$}
\task{$\big< -\!\infty,\;- 2\big]  \cup \big[ 0, 1\big>$}
\task{$\big< -\!\infty,\;- 2\big>  \cup \big<0, 1\big>$}
\task{$\big< -\!\infty,\;- 2\big>  \cup \big<0, 1\big]$}
\task{$\big< -\!\infty,\;- 2\big>  \cup \big[0, 1\big]$}
\end{tasks}
}
\item{En las siguientes proposiciones, escribir (\textrm{V}) si es verdadera o (\textrm{F}) si es falsa.
\begin{enumerate}[label=\roman*.]
\item{El punto medio del segmento de recta de extremos $A = (a,\,b)$ y \\ $B = (b,\,c)$ es $\Bigg[\dfrac{a + b}{2},\,\dfrac{a + c}{2}\Bigg]$.}
\item{Si \quad$M = \Bigg[\dfrac{a + b}{2},\,\dfrac{b + c}{2}\Bigg]$ es el punto medio del segmento de recta donde $A = (a,\,b)$ es uno de los extremos, entonces al otro extremo es $B = (b,\,c)$.}
\item{El punto medio del segmento de recta de extremos $A = (a,\,b)$ y \\ $B = (c,\,d)$ es $\dfrac1{2}\big(a + d,\,b + c\big)$ }
\end{enumerate}
La secuencia correcta, es :
\begin{tasks}(5)
\task{FVV}\task{FVF}\task{FFV}\task{FFF}\task{VFV}
\end{tasks}
}
\item{En las siguientes proposiciones, escribir (\textrm{V}) si es verdadera o (\textrm{F}) si es falsa.
\begin{enumerate}[label=\roman*.]
\item{La recta \;$\mathbb{L} : y = mx + b$\; pasa por el segundo cuadrante, si $m > 0$ y $b < 0$}
\item{La recta $\mathbb{L} : y = mx + b$ pasa por los cuadrantes segundo, tercer y cuarto, si $m < 0$ y $b < 0$}
\item{La recta $\mathbb{L} : y = mx + b$ es creciente, si $m > 0$.}
\end{enumerate}
La secuencia correcta es :
\begin{tasks}(5)
\task{VVF}
\task{VFV}
\task{VFF}
\task{FVV}
\task{FVF}
\end{tasks}
}
\item{La longitud de la circunferencia de ecuación : \\ $(n - 6)\,x^{\,2} + (3n - 8)\,y^{\,2} + (n + 19)\,x - 10y + 10n + 5 = 0$, es : 
\begin{tasks}(5)
\task{$2\sqrt{2}\,\pi$}\task{$3\sqrt{2}\,\pi$}\task{$4\sqrt{2}\,\pi$}\task{$\sqrt{2}\,\pi$}\task{$4\,\pi$}
\end{tasks}
}
\item{Si : $P : (x - a - b)^{\,2} = (a + b)(y - a + b)$ es la ecuación cartesiana de la parábola con vértice en $(4,\,2)$, entonces la ecuación general de dicha parábola es : 
\begin{tasks}(2)
\task{$x^{2} + 8x - 4y + 24 = 0$}
\task{$x^{2} - 8x - 4y + 20 = 0$}
\task{$x^{2} - 8x + 4y + 22 = 0$}
\task{$x^{2} - 8x + 24y + 12 = 0$}
\task{$x^{2} - 8x - 4y + 24 = 0$}
\end{tasks}
}
\item{La ecuación cartesiana de una elipse con eje focal sobre la recta : $\mathbb{L}_{\,1} : x = 2$; lado recto $\;\dfrac{8}{3}\,u\;$, con centro sobre la recta $\mathbb{L}_{\,2} : x - y = 5$ y uno de sus vértices en el punto $(2, 0)$, es :
\begin{tasks}(2)
\task{$\dfrac{(x - 2)^{\,2}}{9} + \dfrac{(y - 3)^{\,2}}{4} = 1$}
\task{$\dfrac{(x - 2)^{\,2}}{4} + \dfrac{(y + 3)^{\,2}}{9} = 1$}
\task{$\dfrac{(x - 2)^{\,2}}{4} + \dfrac{(y - 3)^{\,2}}{9} = 1$}
\task{$\dfrac{(x + 2)^{\,2}}{4} + \dfrac{(y - 3)^{\,2}}{9} = 1$}
\task{$\dfrac{(x - 2)^{\,2}}{8} + \dfrac{(y + 3)^{\,2}}{9} = 1$}
\end{tasks}
}

\end{itemize}
\end{document}